\subsection{I2C Asynchronous Communication}
\label{sec:i2c_async}

The communication begins in the \lstinline[style=myagc]|xfer| method:
\begin{lstlisting}[
    style=myagcs,
    caption={\file{src/drivers/i2c\_async.h}},
    linebackgroundcolor={}
]
template <bool is_read>
int xfer(uint8_t addr,
         uint8_t *cmd, size_t cmd_sz, 
         uint8_t *buf, size_t buf_sz)
{
  // acquire lock before continuing
  mtx.lock(); 

  // ... copy the buffer addresses and buffer sizes

  // initial state is W_CMD
  state.stage = stage_t::W_CMD;
  I2C0.ADDRESS = addr;

  // start transmission and send the first byte
  I2C0.STARTTX = 1;
  I2C0.TXD = *cmd;

  // put the current process to sleep until woken up after completion
  sched::wait(intr_chan);
   
  // ... report any error that might've occured
}
\end{lstlisting}

This method just starts the transmission, and puts the process to sleep. The state
machine is implemented by the interrupt handler:

\begin{lstlisting}[
    style=myagc,
    caption={\file{src/drivers/i2c\_async.h}},
    linebackgroundcolor={\lines{5}{6} \lines{8}{30} \lines{32}{44} \lines{46}{60}
    \lines{62}{67}}
]
void handler(void)
{
  intr_guard guard{I2C0_IRQ};

  if (I2C0.ERROR)
    // ... if there was an error, cancel the communication

  else if (state.stage == stage_t::W_CMD) {
    // ... receive acklowledgement                                      
                                                         +------------+ 
    // still writing the command                         |  W_CMD     | 
    if (state.cmd_idx < state.cmd_sz)                    +------------+ 
      // ... send next byte                                             

    // no data to send or receive, so stop 
    else if (state.data_sz == 0)
      // ... stop

    // if we're writing
    else if (state.mode == mode_t::WRITE) {
      state.stage = stage_t::TX_DATA;
      // ... send the first data byte
    }

    // we're reading
    else {
      state.stage = stage_t::RX_DATA;
      // ... set the SHORTS register and STARTRX = 1
    }
  }

  else if (state.stage == stage_t::TX_DATA) {
    // ... receive acklowledgement
                                                         +------------+              
    // finished writing                                  |  TX_DATA   |               
    if (state.data_idx == state.data_sz)                 +------------+               
      // ... stop                                                       

    // more bytes to write
    else {
      state.stage = stage_t::TX_DATA;
      // ... send the next byte
    }
  }

  else if (state.stage == stage_t::RX_DATA) {
    // ... receive acklowledgement
                                                         +------------+              
    // ... copy byte from I2C0.RXD                       |  RX_DATA   |               
                                                         +------------+               
    // finished reading, move to NACK
    if (state.data_idx == state.data_sz)
      state.stage = stage_t::NACK;

    // continue reading
    else {
      state.stage = stage_t::RX_DATA;
      // ... set the SHORTS register and RESUME = 1
    }
  }

  else if (state.stage == stage_t::NACK) {               
    // ... receive acklowledgement                      +------------+
                                                        |  NACK      | 
    mtx.unlock(); // release the lock                   +------------+
    sched::notify(intr_chan); // then wake up the waiting process
  }
}
\end{lstlisting}

With this implementation, multiple processes can communicate over the I2C bus without
conflicts. All requests are serialized using the \lstinline[style=myagc]|mutex|,
ensuring orderly access to the bus. Additionally, while one process is communicating
over the I2C bus, it does not monopolize the CPU. Instead, it allows other processes to
continue executing, improving overall system efficiency.


